\begin{lemma}{Absch�tzung der unvollst�ndigen Beta--Funktion}
             {AbschaetzungDerUnvollstaendigenBetafunktion}
Seien $n \in \N_{\ge 2}$, $k \in \haken{n-1}$ und $u:[0,\infty[ \rightarrow \R,
t \mapsto\sqrt{\frac 2\pi} \int_t^\infty \exp\left(- \frac{s^2}2 \right)ds$.\\
Dann gilt:
\begin{enumerate}[(1)]
	\item \label{AbschaetzungDerUnvollstaendigenBetafunktion1}
	$\forall\, x \in \!\ ]0, \frac k {n-1}[\, \forall\, m \in \haken{n-k-1}_0:
	\int_0^x t^{n-1-m}(1-t)^m dt \le x^{n-m}(1-x)^m$.
	\item \label{AbschaetzungDerUnvollstaendigenBetafunktion2}
	$\forall\, y \in \R_{>0}\, \forall\, l \in \haken {n-1}_0:
	\int_{y}^\infty tu(t)^{n-l}(1-u(t))^ldt \le
	\int_0^{u(y)}t^{n-l-1}(1-t)^ldt$.
\end{enumerate}
\end{lemma}
\begin{proof}
(1) Seien $x \in\!\  ]0, \frac k {n-1}[$ und $m \in \haken{n-k-1}_0$.\\
Falls $m = 0$, so ist
\[\int_0^{x} t^{n-1-m}(1-t)^m dt 
= \int_0^{x} t^{n-1}dt = \frac 1 n x^n \le x^n = x^{n-m}(1-x)^m\]
Falls $m > 0$, so betrachte
$\alpha: [0, x] \rightarrow \R, t \mapsto t^{n-1-m}(1-t)^m$.\\
Dann ist $\alpha$ offenbar stetig differenzierbar auf $]0,x[$ und es gilt
f�r alle $s \in\!\  ]0, x[:$
\begin{align*}
\alpha'(s)& = (n-1-m)s^{n-m-2}(1-s)^m - ms^{n-m-1}(1-s)^{m-1}\\
&= s^{n-m-2}(1-s)^{m-1}\bigg((n-1-m)(1-s)-ms \bigg)\text.
\end{align*}
Ferner gilt f�r alle $s \in \!\ ]0,x[$:
\[\frac{n-m-1}{n-1} \ge \frac{n-(n-k-1)-1}{n-1} = \frac k{n-1} > x > s\text,\]
also
\[\frac{1-s}s = \frac 1 s - 1 > \frac{n-1}{n-m-1}-1 =
\frac{n-1-(n-m-1)}{n-m-1}= \frac m {n-m-1}\]
und damit $(n-1-m)(1-s) > ms$.\\
Somit ist $\alpha'(s)>0$ f�r alle $s \in\!\  ]0,x[$, also
$\restrict{\alpha}{\!]0,x[}$ monoton wachsend.\\
Damit gilt wegen $\alpha(0) = 0$:
\begin{align*}\int_0^{x} t^{n-1-m}(1-t)^m dt
= \int_0^x \alpha(t)dt \le x  \sup_{t \in [0,x]}\alpha(t)
= x  \alpha(x) = x^{n-m}(1-x)^m\text.\end{align*}
(2) Seien $y \in \R_{>0}$ und $l \in \haken {n-1}_0$.\\
Dann ist $n-l \ge 1$ und es gilt:
\begin{align*}
&\int_y^\infty t u(t)^{n-l}(1-u(t))^ldt
= \int_y^\infty u(t)^{n-l-1}(1-u(t))^l t u(t)dt\\
\overset{\refclaim{EigenschaftenDerKomplementaerenVerteilungsfunktionVonN01}
{EigenschaftenDerKomplementaerenVerteilungsfunktionVonN014}}\le
&\, \int_y^\infty  u(t)^{n-l-1}(1-u(t))^l t \cdot 
\sqrt{\frac 2 \pi} \frac 1 t \exp\left(- \frac 1 2 t^2 \right)dt\\
=&\, \int_y^\infty  u(t)^{n-l-1}(1-u(t))^l \cdot \sqrt{\frac 2 \pi}
\exp\left(- \frac 1 2 t^2 \right)dt\\
\overset{\refclaim{EigenschaftenDerKomplementaerenVerteilungsfunktionVonN01}
{EigenschaftenDerKomplementaerenVerteilungsfunktionVonN012}}=&\,
\int_y^\infty  u(t)^{n-l-1}(1-u(t))^l \abs{u'(t)}dt
= \int_{]y, \infty[}u(t)^{n-l-1}(1-u(t))^l \abs{u'(t)}dt\\
\overset{\substack{\text{Transf.}\\\text{mit }u}}=&\,
\int_{]0,u(y)[}t^{n-l-1}(1-t)^ldt = \int_0^{u(y)} t^{n-l-1}(1-t)^ldt
\end{align*}
\end{proof}